% 工尺谱印本复刻 — 鹊桥仙·shengshi-chengfu
% 参考: 示例/工尺谱/ref-2-upright.png
%
% 复刻范围 (基于 ref 可读部分):
%   鵲橋仙 一名廣寒秋  (调名 + 别名 - 内联子标题)
%   聖世呈符 句  星垣輝朗鴻 韻
%   上 韻  披香侍宴
%   金蓮花炬照回廓 韻
%   上林遊賞 韻    正院 ...
%   金雞叫
%
\documentclass{ltc-guji}

\setmainfont{TW-Kai}

% 收紧版面以匹配 ref-2 印本密度 (默认 21 字/列过松)
\contentSetup{
    n-char-per-col = 12,
    n-column = 8,
}

\begin{document}
\begin{正文}
% 内联子标题 (印本风格: 调名 + 别名独占一列, ref-2 同款)
鵲橋仙\空格\调名注{一名廣寒秋}\换行
% --- 第一句: 聖世呈符·句 / 星垣輝朗鴻·韻 ---
\工尺{聖}{\音[板]{六}\音[眼]{凡}\音[板]{尺}}%
\工尺{世}{\音[眼]{凡}\音[板]{六}}%
\工尺{呈}{\音[眼]{尺}}%
\工尺{符}{\音[板]{凡}}%
\句%
\工尺{星}{\音[板]{六}}%
\工尺{垣}{\音[眼]{凡}}%
\工尺{輝}{\音[板]{六}}%
\工尺{朗}{\音[眼]{尺}\音[板]{上}}%
\工尺{鴻}{\音[板]{凡}}%
\韻%
% --- 第二句: 上 (单字) · 韻 ---
\工尺{上}{\音[板]{尺}}%
\韻%
% --- 第三句: 披香侍宴 ---
\工尺{披}{\音[板]{工}}%
\工尺{香}{\音[眼]{工}}%
\工尺{侍}{\音[板]{六}}%
\工尺{宴}{\音[眼]{尺}}%
% --- 第四句: 金蓮花炬照回廓·韻 ---
\工尺{金}{\音[板]{五}\音[眼]{六}}%
\工尺{蓮}{\音[板]{五}}%
\工尺{花}{\音[眼]{六}}%
\工尺{炬}{\音[板]{五}}%
\工尺{照}{\音[眼]{尺}\音[板]{上}}%
\工尺{回}{\音[眼]{尺}}%
\工尺{廓}{\音[板]{凡}}%
\韻%
% --- 第五句: 上林遊賞·韻 ---
\工尺{上}{\音[板]{六}}%
\工尺{林}{\音[眼]{凡}}%
\工尺{遊}{\音[板]{尺}\音[眼]{上}}%
\工尺{賞}{\音[眼]{凡}}%
\韻%
% --- 第六句: 正院 (片段) ---
\工尺{正}{\音[板]{六}}%
\工尺{院}{\音[眼]{六}}%
% --- 第七句: 金雞叫 ---
\工尺{金}{\音[板]{凡}}%
\工尺{雞}{\音[眼]{六}}%
\工尺{叫}{\音[板]{尺}}
\end{正文}
\end{document}
